\corrigesChapitre{3}{États et changements d’état : le modèle}
\par\noindent\begin{minipage}{\linewidth}
\begin{corrige}{1}{Nommer un passage}
\begin{questions}
\item La fusion.
\item La liquéfaction ; on rencontre aussi le mot condensation liquide.
\end{questions}
\end{corrige}
\end{minipage}
\par\noindent\begin{minipage}{\linewidth}
\begin{corrige}{2}{À l’échelle des particules}
\begin{questions}
\item À l’état gazeux, dans les conditions usuelles comparées.
\item Non. Il épouse le récipient, mais conserve approximativement son volume.
\end{questions}
\end{corrige}
\end{minipage}
\par\noindent\begin{minipage}{\linewidth}
\begin{corrige}{3}{Buée et vapeur}
\begin{questions}
\item Non : ce sont de fines gouttes d’eau liquide.
\item Une partie de la vapeur se liquéfie dans l’air plus froid et forme des gouttelettes qui diffusent la lumière.
\end{questions}
\end{corrige}
\end{minipage}
\par\noindent\begin{minipage}{\linewidth}
\begin{corrige}{4}{Une masse conservée}
\begin{questions}
\item $80\unite g$ d’eau.
\item Non : l’eau liquide occupe moins de volume que la même masse de glace.
\end{questions}
\end{corrige}
\end{minipage}
\par\noindent\begin{minipage}{\linewidth}
\begin{corrige}{5}{Évaporation}
\begin{questions}
\item Une vaporisation par évaporation à la surface.
\item De l’eau, sous forme gazeuse ; les molécules d’eau restent des molécules d’eau.
\end{questions}
\end{corrige}
\end{minipage}
\par\noindent\begin{minipage}{\linewidth}
\begin{corrige}{6}{Lire un palier}
\begin{questions}
\item Solide et liquide.
\item Oui : un transfert d’énergie peut provoquer la fusion sans hausse de température pendant le palier.
\end{questions}
\par\smallskip\noindent\textbf{Pour comprendre.}\enspace Le palier ne signifie pas arrêt du transfert d’énergie : l’énergie reçue permet la fusion pendant la coexistence solide-liquide. La pression et la précision de la mesure font partie de l’interprétation.\par
\pcFigureEx{figures/cor-fusion.pdf}{Au palier, deux états coexistent : le chauffage peut continuer sans hausse de température.}
\end{corrige}
\end{minipage}
\par\noindent\begin{minipage}{\linewidth}
\begin{corrige}{7}{Choisir un état}
\begin{questions}
\item Solide, dans le modèle d’équilibre usuel.
\item Liquide ; on suppose les conditions ordinaires de pression et l’absence de situation particulière.
\end{questions}
\end{corrige}
\end{minipage}
\par\noindent\begin{minipage}{\linewidth}
\begin{corrige}{8}{Reconnaître les étapes}
\begin{questions}
\item D’abord le liquide se réchauffe ; puis la vaporisation se poursuit à température presque constante.
\item La nature et la pureté de l’échantillon, ainsi que la pression de l’expérience.
\end{questions}
\par\smallskip\noindent\textbf{Pour comprendre.}\enspace Lire successivement les portions de la courbe : température croissante, puis constante. Un palier vers 100 °C est compatible avec l’ébullition de l’eau pure à pression usuelle ; la valeur seule ne suffit pas sans les conditions.\par
\pcFigureEx{figures/cor-ebullition.pdf}{La première partie correspond au chauffage du liquide ; le palier à la vaporisation dans les conditions annoncées.}
\end{corrige}
\end{minipage}
\par\noindent\begin{minipage}{\linewidth}
\begin{corrige}{9}{Un bilan de masse}
\begin{questions}
\item $150-138=12\unite g$.
\item Non : une vaporisation suffit à l’expliquer. Il faut considérer l’eau passée dans l’air pour le bilan complet.
\end{questions}
\end{corrige}
\end{minipage}
\par\noindent\begin{minipage}{\linewidth}
\begin{corrige}{10}{Comparer deux courbes}
\begin{questions}
\item A est le plus compatible avec le modèle d’un corps pur.
\item Non : la mesure a une précision limitée et d’autres informations sont nécessaires pour identifier exactement l’espèce.
\end{questions}
\pcFigureEx{figures/cor-deux-courbes.pdf}{A est compatible avec un corps pur ; une courbe seule ne prouve pas l’identité et la pureté exactes.}
\end{corrige}
\end{minipage}
\par\noindent\begin{minipage}{\linewidth}
\begin{corrige}{11}{Ne pas enfermer un chauffage}
\begin{questions}
\item La pression peut augmenter dangereusement ; le protocole doit être défini par le professeur avec du matériel adapté.
\item La vapeur occuperait un grand volume si elle était libre. Si l’espace est imposé, la pression peut fortement croître.
\end{questions}
\end{corrige}
\end{minipage}
\par\noindent\begin{minipage}{\linewidth}
\begin{corrige}{12}{La banquise et le glacier}
\begin{questions}
\item La première flotte déjà dans la mer ; la seconde repose sur un continent.
\item L’eau du glacier s’ajoute à l’océan. Pour la glace flottante, l’effet direct est beaucoup plus faible dans le modèle simple ; elle déplaçait déjà de l’eau.
\end{questions}
\pcFigureEx{figures/cor-glaces.pdf}{La glace continentale apporte de l’eau à l’océan lorsqu’elle fond et s’y écoule ; la banquise flottait déjà.}
\end{corrige}
\end{minipage}
\par\noindent\begin{minipage}{\linewidth}
\begin{corrige}{13}{Critiquer un dessin}
\begin{questions}
\item Conserver les molécules avec la même taille schématique.
\item Leur agitation et leur espacement moyen peuvent changer ; on peut représenter cela avec des flèches et un écartement adapté.
\end{questions}
\pcFigureEx{figures/cor-particules-chaudes.pdf}{Le modèle conserve les entités ; le chauffage modifie notamment leur agitation, sans les transformer en grosses molécules.}
\end{corrige}
\end{minipage}
\par\noindent\begin{minipage}{\linewidth}
\begin{corrige}{14}{Pression et température}
\begin{questions}
\item Non : la température d’ébullition dépend de la pression.
\item Il faut comparer les pressions, ainsi que la pureté de l’eau. La référence de 100 °C suppose une pression voisine de 1,013 bar.
\end{questions}
\end{corrige}
\end{minipage}
\par\noindent\begin{minipage}{\linewidth}
\begin{corrige}{15}{Construire un argument}
\begin{questions}
\item Non, la même température peut correspondre à des situations différentes.
\item Observer la présence de glace, ou suivre l’évolution pendant un transfert d’énergie connu. La température seule ne décrit pas la proportion des états.
\end{questions}
\end{corrige}
\end{minipage}
\par\noindent\begin{minipage}{\linewidth}
\phantomsection\label{cor-defi-3}\pcSolutionDefi{3}{Le flacon qui semble perdre de la matière}\begin{questions}
\item $312-300=12\unite g$.
\item De l’eau peut s’être vaporisée et avoir quitté la partie pesée : le bilan doit inclure cette eau.
\item Le professeur peut utiliser un montage adapté pour condenser et recueillir la vapeur, puis comparer l’ensemble des masses. Ne jamais improviser un chauffage hermétique.
\end{questions}
\end{minipage}
\par\noindent\begin{minipage}{\linewidth}
\phantomsection\label{cor-activite-3}\pcSousTitre{Pistes pour l’investigation}
Pendant la coexistence eau-glace, la température reste proche de 0 °C sous pression usuelle, malgré un transfert d’énergie. Les petites variations doivent être comparées à la précision de la sonde. La fonte d’une glace continentale ajoute de l’eau à la mer ; la glace déjà flottante déplaçait auparavant une quantité d’eau liée à sa masse. La dilatation de l’océan est un autre mécanisme.\par\end{minipage}
